\chapter{Conclusion}
\section{Key Findings}
This project implemented and evaluated DE/rand/1/bin for RBF-SVR hyperparameter optimization on concrete compressive strength. A fixed group-disjoint protocol, fold-local scaling and equal initialization-inclusive search budgets allowed a controlled comparison with default SVR and Random Search.

DE achieves mean holdout RMSE of $6.2250\pm0.1104$ MPa, MAE of $4.3143\pm0.0526$ MPa and $R^2$ of $0.8324\pm0.0060$. The mean RMSE reduction from default SVR is approximately 37.2\%. Random Search achieves essentially the same holdout accuracy. The experiment therefore demonstrates the value of tuning and an auditable evolutionary implementation, while providing no basis for claiming significant general superiority of DE.
A pre-registered extension found that the input representation limited the frozen results more than the optimizer. With $\ln(\text{age})$ for all three methods, holdout RMSE is 7.4780 MPa for default SVR, $4.9700\pm0.2656$ MPa for Random Search and $4.8402\pm0.1039$ MPa for DE. DE is slightly better and less variable than Random Search here, but with five seeds on one partition this is an observation, not a significant difference. DE variants (current-to-best/1/bin, jDE) perform alike.
\section{Key Learnings}
Data isolation and honest experiment accounting are central to a useful optimizer comparison. Repeated input tuples need group-aware splits; scaling must be fitted within folds; every candidate evaluation must count toward the budget. The demonstration model and five-seed aggregate answer different questions and must remain distinct. Examining the input representation was more productive than examining the optimizer, and disclosing how the log-age idea arose keeps that result honest.

Preserving original provenance, complete histories and fixed selections made it possible to audit results across environments. This reproducibility is a stronger project outcome than an unsupported accuracy claim based on a favorable seed.
